\documentclass[11pt]{article}
\usepackage{preamble}
\renewcommand{\answer}[1]{}

\begin{document}

\ladderheading{AP Statistics \textperiodcentered\ Unit 3 \textperiodcentered\ Topic 3.3 \textperiodcentered\ B: 2026-11-24 \textperiodcentered\ E: 2026-12-23}{Six Activations in Twelve Users}

\focusbanner{TODAY'S PROBLEM}

Suppose a developer's model says $p_0=0.20$ of new app users activate voice commands. An SRS of 12 users includes 6 activations, so $\hat p=0.50$. The chart shows 200 simulated samples of 12 independent users under this model.\par\smallskip \textbf{Mara:} ``A normal model for $\hat p$ gives $P(\hat p\ge 0.50)\approx0.005$, so that is the chance probability I would report.''\par \textbf{Jules:} ``The simulated distribution is not very normal. I would estimate the chance probability directly from its tail.''\par
\begin{center}
\begin{tikzpicture}
\begin{axis}[
  width=0.98\linewidth,height=1.75in,
  ybar,bar width=10pt,xmin=-0.6,xmax=6.6,ymin=0,ymax=70,
  xtick={0,1,2,3,4,5,6},ytick={0,20,40,60},
  xlabel={Activations in a sample of 12},ylabel={Simulated samples},
  tick label style={font=\small},label style={font=\small},
  nodes near coords,nodes near coords style={font=\scriptsize},
  axis lines*=left
]
\addplot[draw=dayblue,fill=dayblue!20] coordinates {(0,14) (1,40) (2,58) (3,47) (4,25) (5,12) (6,4)};
\end{axis}
\end{tikzpicture}
\end{center}
\begin{firsttakebox}{FIRST TAKE --- before the video (5 min)}
Whose method would you trust more, Mara's smooth curve or Jules's repeated trials? Give one reason from the picture.
\workspace{0.9in}
\end{firsttakebox}

\begin{callout}{\IconBook\ CHECK THE SHAPE BEFORE USING NORMAL (keep on the page)}{calloutgreen}
Under a proposed proportion $p_0$, check expected counts $np_0\geq10$ and $n(1-p_0)\geq10$ before using a normal approximation.
A matched simulation estimates the chance probability by the number of simulated results at least as large as the observed result, divided by the total simulations.
For this task, a chance probability below $5\%$ counts as convincing evidence against the proposed model.
Evidence against a model does not prove a particular cause.
\end{callout}

\newpage
\textbf{Finish after the video:}\par\smallskip

\dokbadge{1}~\textbf{(a)} State the most common number of activations in the simulated samples and describe the distribution as roughly symmetric, skewed left, or skewed right.\par
\workspace{0.7in}

\dokbadge{2}~\textbf{(b)} Use the chart to estimate $P(\hat p\ge0.50)$ under $p_0=0.20$. Then calculate $np_0$ and $n(1-p_0)$ and assess the expected-count condition for a normal approximation.\par
\workspace{1.4in}

\dokbadge{3}~\textbf{(c)} In 3--4 sentences, choose a method using the simulated tail and either the shape or expected counts. Apply the $5\%$ evidence rule. Explain how reporting 0.005 instead of 0.020 would change a reader's impression.\par
\workspace{2.0in}

\begin{sentenceframebox}\raggedright\textbf{Frame:} I would use \blank[1.2in] because the expected counts are \blank[1.5in] and the simulated shape is \blank[1.4in].\end{sentenceframebox}

\begin{sentenceframebox}\raggedright\textbf{Frame:} The tail estimate \blank[0.8in] means \blank[2.3in]; reporting $0.005$ instead would make a reader think \blank[2.2in].\end{sentenceframebox}

\begin{summaryexitbox}{EXIT --- turn this in}
One thing the video changed about my first take: \hrulefill\par
\workspace{0.4in}
\end{summaryexitbox}

\end{document}
